\subsection{布绕尔群的定义}

我们用 \(\mathrm{Br}(\mathrm{K})\) 表示 \(M_{K}\) 中形如 {[}A{]} 的元素所成之集，其中 A 是 K 上一个有限次数的中心单代数。

命题 2. —— 集合 Br(K) 是 \(\mathcal{M}_{\mathrm{K}}\) 的可逆元素所成之集，其中 Br(K) 的元素 {[}A{]} 的逆是 {[}A°{]}。因此，幺半群 \(M_{K}\) 上的合成律给 Br(K) 装备了一个交换群结构。

设 A 是 K 上一个有限次数的中心单代数。代数 \(A^{o}\) 是 K 上有限次数的中心单代数（VIII, 第 251 页）。代数 \(A \otimes_{K} A^{o}\) 同构于一个矩阵代数 \(\mathbf{M}_{n}(K)\)，其中 \(n^{2} = \dim(A)\)（VIII, 第 252 页，定理 1）。因此我们有 \([A][A^{o}] = [A \otimes_{K} A^{o}] = [K]\)（引理 1），这就证明 {[}A{]} 是可逆的，其逆为 \([A^{o}]\)。

反之，设 A 是一个有限次数的 K-代数。若 {[}A{]} 在 \(M_{K}\) 中是可逆的，则存在一个有限次数的 K-代数 B，使得 \([A][B]=[K]\)；由公式 (2) 与引理 1，这意味着 K-代数 \(A \otimes_{K} B\) 同构于一个矩阵代数 \(\mathbf{M}_{n}(\mathrm{K})\)，其中 \(n \geqslant 1\)。由 VIII, 第 251 页的注 1，代数 A 于是是中心单的。

定义 1. —— 交换群 Br(K) 称为域 K 的布绕尔群。

引理 2. —— 设 I 与 J 是有限集，k 是一个交换环，A 与 B 是 k-代数。用 \(\mathbf{M}_{\mathrm{I}}(\mathrm{A})\) 表示元素在 A 中的 (I,I) 型方阵所成的 k-代数，并类似地定义 k-代数 \(\mathbf{M}_{\mathrm{J}}(\mathrm{B})\) 与 \(\mathbf{M}_{\mathrm{I}\times\mathrm{J}}(\mathrm{A}\otimes_{\mathrm{K}}\mathrm{B})\)。存在唯一的 k-代数同构

\[
\varphi : \mathbf {M} _ {\mathrm{I}} (\mathrm{A}) \otimes_ {k} \mathbf {M} _ {\mathrm{J}} (\mathrm{B}) \longrightarrow \mathbf {M} _ {\mathrm{I} \times \mathrm{J}} (\mathrm{A} \otimes_ {k} \mathrm{B})
\]

使得 \(\varphi\big((a_{ii'})\otimes(b_{jj'})\big)\) 是指标为 \(((i,j),(i',j'))\) 的元素等于 \(a_{ii'}\otimes b_{jj'}\) 的矩阵。

具有引理所述性质的 k-线性双射 \(\varphi\) 的存在性由张量积与直和的相容性得出（II, §3, No.~7, 第 255 页，命题 7）。\(\varphi\) 是代数同构这一事实由乘积矩阵的定义得出。

命题 3. —— 设 A 与 B 是有限次数的中心单 K-代数。下列性质等价：

\begin{enumerate}
\def\labelenumi{(\roman{enumi})}
\item
  我们在布绕尔群 Br(K) 中有 {[}A{]} = {[}B{]}。
\item
  存在一个整数 \(t \geqslant 1\)，使得 K-代数 \(\mathrm{A} \otimes_{\mathrm{K}} \mathrm{B}^{\circ}\) 同构于矩阵代数 \(\mathbf{M}_t(\mathrm{K})\)。
\item
  存在严格正的整数 \(r\) 与 \(s\)，使得 K-代数 \(\mathrm{A} \otimes_{\mathrm{K}} \mathbf{M}_r(\mathrm{K})\) 与 \(\mathrm{B} \otimes_{\mathrm{K}} \mathbf{M}_s(\mathrm{K})\) 同构。
\item
  存在一个包含 K 的体 D 以及整数 \(m \geqslant 1\) 与 \(n \geqslant 1\)，使得 A 同构于 \(\mathbf{M}_{m}(\mathrm{D})\) 且 B 同构于 \(\mathbf{M}_{n}(\mathrm{D})\)。
\item
  K-代数 A 与 B 是森田等价的。
\end{enumerate}

假设我们有 \([A]=[B]\)。由于 \([B^{o}]\) 是 {[}B{]} 在布绕尔群中的逆，我们有 \([K]=[B][B^{o}]=[A][B^{o}]=[A\otimes_{K}B^{o}]\)。由引理 1，存在一个整数 \(t\geqslant1\)，使得代数 \(A\otimes_{K}B^{o}\) 与 \(\mathbf{M}_{t}(K)\) 同构。因此 (i) 蕴含 (ii)。

假设 (ii) 成立。由于 \(B^{o} \otimes_{K} B\) 同构于一个矩阵代数 \(\mathbf{M}_{s}(K)\)，其中 \(s \geqslant 1\)（VIII, 第 252 页，定理 1），代数 \(A \otimes_{K} B^{o} \otimes_{K} B\) 一方面同构于 \(A \otimes_{K} M_{s}(K)\)，另一方面同构于 \(\mathbf{M}_{t}(K) \otimes_{K} B\)。因此性质 (ii) 蕴含性质 (iii)。

假设 (iii) 成立。由韦德伯恩定理（VIII, 第 120 页，定理 1），存在整数 \(m \geqslant 1\) 与 \(n \geqslant 1\) 以及 K 上有限次数且以 K 为中心的体 D 与 D'，使得 A 同构于 \(\mathbf{M}_m(\mathrm{D})\) 且 B 同构于 \(\mathbf{M}_n(\mathrm{D}')\)。于是代数 \(\mathrm{A} \otimes_{\mathrm{K}} \mathbf{M}_r(\mathrm{K})\) 同构于 \(\mathbf{M}_{mr}(\mathrm{D})\)，而代数 \(\mathrm{B} \otimes_{\mathrm{K}} \mathbf{M}_{r'}(\mathrm{K})\) 同构于 \(\mathbf{M}_{nr'}(\mathrm{D}')\)（引理 2）。由 VIII, 第 121 页的推论 2，K-代数 D 与 D' 同构。因此 (iii) 蕴含 (iv)。

假设 (iv) 成立。代数 A 与 D 是森田等价的，代数 B 与 D 亦然（VIII, 第 114 页，例 2）。因此 A 与 B 是森田等价的，故 (iv) 蕴含 (v)。

最后，蕴含关系 \((v)\Rightarrow(i)\) 就是集合 \(\mathrm{Br}(K)\) 的定义。

当命题的这些等价的性质成立时，我们说代数 A 与 B 是相似的。

推论. —— 设 A 与 B 是 K 上有限次数的中心单代数。则 A 与 B 同构当且仅当它们相似且次数相同。

这由命题 3 的性质 (i) 与 (iv) 的等价性以及 \(\dim_{\mathrm{K}}(\mathbf{M}_n(\mathrm{D})) = \dim_{\mathrm{K}}(\mathrm{D}) \times n^2\) 这一事实得出。

命题 4. —— 设 \(\mathcal{H}_{\mathrm{K}}\) 是作为体的有限次数中心 K-代数的类所成之集。映射 \(\mathrm{D} \mapsto [\mathrm{D}]\)（从 \(\mathcal{H}_{\mathrm{K}}\) 到 \(\mathrm{Br}(\mathrm{K})\) 的）是双射的。

这由 VIII, 第 278 页的引理 1 与韦德伯恩定理（VIII, 第 120 页，定理 1）得出。

